\section{Introduction}
\label{sec:intro}

Infrared and quantum-cascade-laser histology classify tissue pixel by
pixel from a spectrum of several hundred channels, and one family of
models used for this task is serial: a small encoder compresses each
pixel's spectrum, and a much larger spatial network reads the encoded
neighbourhood \citep{berisha2019deep,oleary2026spatial}. Practitioners
often pretrain or freeze the spectral stage, and a companion tokenization
study (anonymous, under review) reports a frozen random encoder competing
with a jointly trained one. Such comparisons do not settle whether joint
training starves the encoder: the spatial stage can predict the label
from the neighbourhood, and a model that learns to do so may have little
incentive to improve access to the target's centre-spectrum cue.

This paper asks how the relative training speed of the two modules
affects what the encoder learns, keeping four easily conflated
measurements apart: adaptation of the encoder, linear accessibility of
the spectral cue in its output, reliance of the trained classifier on
context, and what a specified retraining procedure can recover afterwards.
We prove one instance exactly: in a linear-encoder model with a
replicated contextual readout, a faster contextual pathway suppresses the
update of the spectral coefficient at matched fit, the fitted model fails under contextual
reversal under explicit finite-width conditions, and readout normalization
removes the width dependence (Section~\ref{sec:theorem}). We then measure,
with interventions on training speed, how each responds in a nonlinear
convolutional model on synthetic data (Section~\ref{sec:synthetic}) and
with measured tissue spectra in constructed neighbourhoods
(Section~\ref{sec:real}), each against a control that removes the
proposed cause.

\paragraph{Contributions.}
\begin{itemize}
\item \textbf{Theorem~\ref{thm:serial}.} For the serial cross-entropy model
  with $M$ replicated contextual readouts, gradient flow from general
  initialization moves the spectral coefficient at the matched fitting
  margin by at most a fraction $1/(1+Mv_0^2)$ of the update a spectral-only
  model needs; the fitted model misclassifies every example under
  contextual reversal when the initial spectral coefficient is below half
  the margin and $Mv_0^2>m/(m-2a_0)$, with expected error tending to
  $\Phi(m/2\sigma)$ over isotropic Gaussian initialization as $M\to\infty$;
  and a $M^{-1/2}$-normalized readout removes the width dependence, though
  it need not remove the failure.
\item \textbf{Synthetic intervention.} In a ReLU convolutional head over a
  linear encoder, with two cues of matched signal-to-noise ratio, an
  informative context suppresses the encoder's spectral learning at
  matched fit at every width. Slowing the whole head by $256\times$ raises
  the spectral cue's linear accessibility from $0.66$ to $0.85$,
  monotonically in every seed, while the trained classifier's reversal
  accuracy moves from $0.12$ to $0.14$; a normalized readout reduces the
  suppression at large width; an uninformative context removes it.
  Retraining a fresh head on decorrelated context recovers shifted
  accuracy of $0.75$--$0.84$ after the slow head and $0.54$--$0.62$ after
  the fast one, near the random-encoder baseline ($0.53$--$0.60$); a
  two-layer ReLU encoder reproduces the ordering.
\item \textbf{Real spectra.} With measured centre spectra in constructed
  neighbourhoods, informative-context arms are strongly context-reliant
  in a high- and a low-accessibility regime; in the latter the encoder's
  adaptation is limited relative to the uninformative-context control,
  with a rate response that is larger at the smaller constructed
  amplitude. Where the spectral cue is accessible, retraining
  the head recovers $0.90$--$0.96$ shifted accuracy from every encoder, a
  random one included: the original failure coexisted with spectral
  information that the specified retraining could use. On a public
  hyperspectral scene the constructed-context experiment reproduces
  suppression and reliance, with rate and recovery effects inconsistent
  across seeds; in natural $3\times3$ neighbourhoods, on tissue and on
  that scene, contextual reliance occurs without a detected
  encoder-accessibility deficit (Appendices~\ref{app:exp4}
  and~\ref{app:pavia}).
\item \textbf{Scope.} Scalar curvature on a production model does not
  identify the mechanism, and an earlier apparent freezing advantage was
  protocol-confounded (Appendix~\ref{app:scope}); no mitigation is tested
  on natural context.
\end{itemize}

\paragraph{Related work.}
That one predictive feature can suppress the learning of another is
gradient starvation \citep{pezeshki2021gradient}, and a preference for
simpler features that fails under shift is simplicity bias
\citep{shah2020pitfalls}; modality competition in multimodal training
\citep{wang2020what,peng2022balanced,huang2022modality,du2023unimodal}
studies the same imbalance between input streams. That core features can
be present in a representation while a spurious feature is used, and that
last-layer retraining can then restore robustness, is the observation of
\citet{kirichenko2023dfr}; we test how the original relative module rate
changes recovery by retraining the entire nonlinear head from a frozen
encoder, comparing original encoders at matched fit against a control.
In infrared histology, retained performance under a learned
sixteen-feature bottleneck, in spatial and in spectrum-only classifiers,
has been read as showing that tissue classification needs few spectral
features \citep{oleary2026spatial}, and CNN classifiers are reported
insensitive to the compression method \citep{mueller2023dimred}.
Reliance is a different question: a fitted predictor's spatial reliance
does not show useful spectral information to be dispensable, since at one
bottleneck dimension our rate interventions give similarly
context-reliant predictors whose frozen encoders support very different
shifted accuracy after the same retraining; a learned twelve-dimensional
bottleneck within tolerance of none while sixteen principal components
lose performance (Appendix~\ref{app:exp4}) shows dependence on the
representation, not a need for high-dimensional encoder output. The mathematics
uses the balance invariant of gradient flow in layered linear models
\citep{astra_du2018balance,astra_saxe2014,astra_yun2021} next to
lazy-training analyses \citep{chizat2019lazy,yang2021tensor}; we measure
learning outcomes rather than curvature, whose observed width trend in
the production model is sensitive to normalization
\citep{karakida2019universal,vanlaarhoven2017l2}
(Appendix~\ref{app:scope}); freezing in practice
\citep{zhang2022layers,coil2025freezing} is the setting, not a remedy we
evaluate.
